% Pala con profilo NACA 2412 approssimato.
% Uso: \DrawBlade{x}{y}{angolo}

\newcommand{\DrawBlade}[3]{%
  \begin{scope}[shift={(#1,#2)},rotate=#3,xscale=2.2,yscale=2.2]

    % Profilo della pala di turbina (curvatura a C)
    \draw[thick,fill=gray!10,line join=round]
      % Estradosso (suzione - molto convesso)
      plot[smooth] coordinates {
        (0.000, 0.000)
        (0.050, 0.120)
        (0.150, 0.220)
        (0.300, 0.280)
        (0.450, 0.270)
        (0.600, 0.220)
        (0.750, 0.140)
        (0.900, 0.060)
        (1.000, 0.000)
      }
      --
      % Intradosso (pressione - nettamente concavo)
      plot[smooth] coordinates {
        (1.000, 0.000)
        (0.850, 0.040)
        (0.700, 0.080)
        (0.550, 0.120)
        (0.400, 0.140)
        (0.250, 0.130)
        (0.120, 0.080)
        (0.050, 0.030)
        (0.000, 0.000)
      }
      -- cycle;

    % Linea media (Camber line)
    \draw[dashed,gray!65,thin]
      plot[smooth] coordinates {
        (0.000, 0.000)
        (0.150, 0.150)
        (0.350, 0.210)
        (0.600, 0.170)
        (0.800, 0.090)
        (1.000, 0.000)
      };

  \end{scope}%
}


% ============================================================
% PARAMETRI MODIFICABILI
% ============================================================
\def\Denom{50}

% Pale
\def\BladeAngle{105}
\def\BladeY{-6.8}

% Triangolo di ingresso
\def\OoneX{0}
\def\OoneY{0}

% Estremo di V1
% VoneX diverso da OoneX rende V1 inclinato.
\def\VoneX{0/\Denom}
\def\VoneY{-154/\Denom}

% Estremo di W1
\def\WoneX{-248/\Denom}
\def\WoneY{-154/\Denom}

% Valori visualizzati
\def\VoneValue{154}
\def\WoneValue{292}
\def\UoneValue{248}

% Angoli visualizzati
\pgfmathsetmacro{\AlphaOneValue}{round(10*atan2(abs(\VoneX),abs(\VoneY)))/10}
\def\BetaOneValue{58.1}

% Triangolo di uscita
\def\OtwoX{9}
\def\OtwoY{\BladeY-1}

% Estremo di V2
\def\VtwoX{101/\Denom+\OtwoX}
\def\VtwoY{-154/\Denom+\OtwoY}

% Estremo di W2
\def\WtwoX{-147/\Denom+\OtwoX}
\def\WtwoY{-154/\Denom+\OtwoY}

% Valori visualizzati
\def\VtwoValue{184}
\def\WtwoValue{213}
\def\UtwoValue{248}

% Angoli visualizzati
\def\AlphaTwoValue{33.3}
\def\BetaTwoValue{43.6}



\begin{figure}[H]
\centering
\begin{tikzpicture}[
  scale=0.8,
  font=\fontsize{7}{8}\selectfont,
  absolute/.style={ultra thick,cyan!80!blue,-{Stealth[length=3mm,width=2.2mm]}},
  relative/.style={ultra thick,orange,-{Stealth[length=3mm,width=2.2mm]}},
  peripheral/.style={ultra thick,green!70!black,-{Stealth[length=3mm,width=2.2mm]}},
  reference/.style={dashed,gray!60,thin},
  velocitylabel/.style={text=black,fill=white,inner sep=2pt},
  angletext/.style={text=black,fill=white,inner sep=1pt}
]

% ============================================================
% TRIANGOLO DI INGRESSO
% ============================================================

\coordinate (O1) at (\OoneX,\OoneY);
\coordinate (V1) at (\VoneX,\VoneY);
\coordinate (W1) at (\WoneX,\WoneY);
\coordinate (VerticalOne) at (\OoneX,-3.1);

% Linea verticale di riferimento
\draw[reference] (O1) -- (VerticalOne);

% Vettore V1
\draw[absolute]
  (O1) -- (V1)
  node[velocitylabel,midway,right=2pt]
  {$v_{1}=\VoneValue\,\mathrm{m/s}$};

% Vettore W1
\draw[relative]
  (O1) -- (W1)
  node[velocitylabel,midway,below left=1pt]
  {$w_{1}=\WoneValue\,\mathrm{m/s}$};

% Vettore U1
\draw[peripheral]
  (W1) -- (V1)
  node[velocitylabel,midway,below=5pt]
  {$U=\UoneValue\,\mathrm{m/s}$};

% Arco beta 1
\pic[
  draw,
  angle radius=1.25cm
]
  {angle=W1--O1--VerticalOne};

\node[angletext]
  at (1,-0.7)
  {$\beta_1=\BetaOneValue^\circ$};

% Arco alpha 1
% Diventa visibile quando V1 viene inclinato.
%\pic[
%  draw,
%  angle radius=0.72cm
%]
%  {angle=VerticalOne--O1--V1};

%\node[angletext]
%  at (1.6,0)
%  {$\alpha_{1.5} = \pgfmathprintnumber[fixed, precision=1]{\AlphaOneValue}^\circ$}


% ============================================================
% PALE NACA
% ============================================================

\DrawBlade{0}{\BladeY}{\BladeAngle}
\DrawBlade{2}{\BladeY}{\BladeAngle}
\DrawBlade{4}{\BladeY}{\BladeAngle}
\DrawBlade{6}{\BladeY}{\BladeAngle}
\DrawBlade{8}{\BladeY}{\BladeAngle}




% ============================================================
% FRECCIA DI ROTAZIONE
% ============================================================

\node[
  single arrow,
  draw=green!70!black,
  fill=white,
  thick,
  minimum height=1cm,
  minimum width=0.5cm,
  single arrow head extend=0.1cm,
  anchor=east,
  shape border rotate=0
]
  (rotor) at (-1,\BladeY+1) {};

\node[
  left=4pt,
  align=center,
  font=\sffamily\scriptsize
]
  at (rotor.west)
  {\textbf{Rotazione}\\\textbf{del rotore}};


% ============================================================
% PIANO DI USCITA
% ============================================================

\draw[gray!65,thin]
  (-1,\BladeY) -- (9,\BladeY);

\node[below,text=gray!65,font=\scriptsize]
  at (3.2,\BladeY)
  {piano di uscita del rotore};


% ============================================================
% TRIANGOLO DI USCITA
% ============================================================

\coordinate (O2) at (\OtwoX,\OtwoY);
\coordinate (V2) at (\VtwoX,\VtwoY);
\coordinate (W2) at (\WtwoX,\WtwoY);
\coordinate (VerticalTwo) at (\OtwoX,\BladeY-4);

% Linea verticale di riferimento
\draw[reference] (O2) -- (VerticalTwo);

% Vettore V2
\draw[absolute]
  (O2) -- (V2)
  node[velocitylabel,midway,below right=2pt]
  {$v_2=\VtwoValue\,\mathrm{m/s}$};

% Vettore W2
\draw[relative]
  (O2) -- (W2)
  node[velocitylabel,midway,left=10pt]
  {$w_2=\WtwoValue\,\mathrm{m/s}$};

% Vettore U2
\draw[peripheral]
  (W2) -- (V2)
  node[velocitylabel,midway,below=8pt]
  {$U=\UtwoValue\,\mathrm{m/s}$};

% Arco alpha 2
\pic[
  draw,
  angle radius=0.6cm
]
  {angle=VerticalTwo--O2--V2};

\node[angletext]
  at (10.4,\BladeY-1.5)
  {$\alpha_2=\AlphaTwoValue^\circ$};

% Arco beta 2
\pic[
  draw,
  angle radius=1cm
]
  {angle=W2--O2--VerticalTwo};

\node[angletext]
  at (8.2,\BladeY-3)
  {$\beta_2=\BetaTwoValue^\circ$};



\draw[->, thick] (9,0) -- (10.5, 0) node[right] {$\theta$};
% Asse theta (verticale)
\draw[->, thick] (9,0) -- (9, -1.5) node[below] {$z$};
% Origine
\fill (0,0) circle (1.5pt);

  
\end{tikzpicture}
    \caption{Triangoli di velocità all'ingresso e all'uscita del rotore del primo stadio del compressore.}
    \label{fig:compressore:triangoli}
\end{figure}

